%% kapa: exam=1 round=37 year=2026 period=1 form=A subject=economics subject_ko=경제학원론 seq=19 number=59
%% source: Q-Net 감정평가사 기출문제 (cst00309) artlSeq=5258924 file=raw/1008/5258924/2026년 제37회 감정평가사 1차 1교시.pdf
%% conversion: kordoc 4.24.51 -> tools/gen1.py (auto); review=proofread-sonnet; flags=none
\begin{question}{59}
소비의 경합성과 배제성을 기준으로 재화를 구분할 때, ( )에 들어갈 내용으로 옳은 것은?
\begin{framedtext}
○ 소비의 경합성이 있으면서 배제성이 없는 재화는 ( ㄱ )이다. \par
○ 소비의 경합성이 없으면서 배제성도 없는 재화는 ( ㄴ )이다. \par
\end{framedtext}
\begin{choices}
  \item ㄱ: 사유재, ㄴ: 공공재
  \item ㄱ: 공유재, ㄴ: 공공재
  \item ㄱ: 공공재, ㄴ: 단체재
  \item ㄱ: 단체재, ㄴ: 공유재
  \item ㄱ: 공유재, ㄴ: 사유재
\end{choices}
\end{question}
\answer{2}
